\chapter{Interfaces de sistemas operacionais}
\label{CH:UNIX}

A função de um sistema operacional é compartilhar um computador entre
múltiplos programas e fornecer um conjunto de serviços mais úteis
do que o suporte oferecido apenas pelo hardware.
Um sistema operacional gerencia e abstrai
o hardware de baixo nível, de modo que, por exemplo,
um processador de texto não precise se preocupar com qual tipo
de disco está sendo utilizado.
Um sistema operacional compartilha o hardware entre vários programas para que
eles executem (ou aparentem executar) ao mesmo tempo.
Por fim, sistemas operacionais fornecem maneiras controladas para que os programas
interajam, de modo que possam compartilhar dados ou trabalhar em conjunto.

Um sistema operacional fornece serviços aos programas de usuário por meio de uma interface.
\index{projeto de interface}
Projetar uma boa interface acaba sendo
uma tarefa difícil. Por um lado, gostaríamos que a interface fosse
simples e enxuta, pois isso facilita a implementação correta.
Por outro lado,
podemos ser tentados a oferecer muitos recursos sofisticados às aplicações.
O truque para resolver essa tensão é projetar interfaces que se baseiem em poucos
mecanismos que possam ser combinados para oferecer grande generalidade.

Este livro utiliza um único sistema operacional como exemplo concreto para
ilustrar conceitos de sistemas operacionais. Esse sistema operacional,
o xv6, fornece as interfaces básicas introduzidas por Ken Thompson e
Dennis Ritchie no sistema Unix~\cite{unix}, além de imitar o projeto interno do Unix.
O Unix fornece uma
interface enxuta cujos mecanismos se combinam bem, oferecendo um surpreendente
grau de generalidade. Essa interface foi tão bem-sucedida que
sistemas operacionais modernos—BSD, Linux, macOS, Solaris e até, em menor grau,
o Microsoft Windows—possuem interfaces semelhantes ao Unix.
Compreender o xv6 é um bom começo para entender qualquer um desses
sistemas e muitos outros.

Como mostra a 
Figura~\ref{fig:os},
o xv6 assume a forma tradicional de um
\indextext{kernel} (núcleo),
um programa especial que fornece
serviços aos programas em execução.
Cada programa em execução, chamado de
\indextext{processo},
possui memória contendo instruções, dados e uma pilha. As
instruções implementam o
cálculo do programa. Os dados são as variáveis sobre as quais
o cálculo atua. A pilha organiza as chamadas de procedimentos do programa.
Um computador geralmente possui muitos processos, mas apenas um único
kernel.

Quando um
processo precisa invocar um serviço do kernel, ele chama
uma \indextext{chamada de sistema} (system call),
uma das chamadas
da interface do sistema operacional.
A chamada de sistema entra no kernel;
o kernel realiza o serviço e retorna.
Assim, um processo alterna entre a execução no
\indextext{espaço do usuário}
e no
\indextext{espaço do kernel}.

Como será descrito em detalhes nos capítulos seguintes, o kernel utiliza os mecanismos de proteção de hardware fornecidos por uma
CPU\footnote{
Este texto geralmente se refere ao elemento de hardware que executa uma
computação com o termo \indextext{CPU}, sigla para unidade central
de processamento (central processing unit). Outras documentações (por exemplo, a especificação RISC-V)
também utilizam os termos processador, núcleo (core) e hart em vez de CPU.
}
para
garantir que cada processo executando no espaço do usuário possa acessar apenas
sua própria memória.
O kernel executa com os privilégios de hardware necessários para
implementar essas proteções; os programas de usuário executam sem
esses privilégios.
Quando um programa de usuário invoca uma chamada de sistema, o hardware
eleva o nível de privilégio e começa a executar uma função pré-configurada
no kernel.

\begin{figure}[t]
\center
\includegraphics[scale=0.5]{fig/os.pdf}
\caption{Um kernel e dois processos de usuário.}
\label{fig:os}
\end{figure}

O conjunto de chamadas de sistema que um kernel fornece
é a interface que os programas de usuário enxergam.
O kernel do xv6 fornece um subconjunto dos serviços e chamadas de sistema
que os kernels Unix tradicionalmente oferecem.  
A Figura~\ref{fig:api} 
lista todas as chamadas de sistema do xv6.

O restante deste capítulo descreve os serviços do xv6—processos, memória,
descritores de arquivos, pipes e sistema de arquivos—e os ilustra com
trechos de código e discussões sobre como o \indextext{shell},
a interface de linha de comando do Unix, os utiliza.
O uso das chamadas de sistema pelo shell ilustra o quanto essas interfaces
foram cuidadosamente projetadas.

O shell é um programa comum que lê comandos do usuário
e os executa.
O fato de o shell ser um programa de usuário, e não parte do kernel,
ilustra o poder da interface de chamadas de sistema: não há nada de
especial sobre o shell.
Isso também significa que o shell pode ser facilmente substituído; como resultado,
sistemas Unix modernos possuem uma variedade de
shells disponíveis, cada um com sua própria interface
e recursos de script.
O shell do xv6 é uma implementação simples da essência do
shell Bourne do Unix. Sua implementação pode ser encontrada em 
\lineref{user/sh.c}{1}.

%% 
%% 	Processes and memory
%% 
\section{Processos e memória}

Um processo no xv6 consiste em memória no espaço do usuário (instruções, dados e pilha)
e estado por processo, privado do kernel.
O xv6
\indextext{compartilha o tempo} (time-share)
entre processos: ele alterna transparentemente as CPUs disponíveis
entre o conjunto de processos que aguardam execução.
Quando um processo não está em execução, o xv6 salva os registradores da CPU do processo,
restaurando-os quando for executá-lo novamente.
O kernel associa um identificador de processo, ou
\indexcode{PID},
a cada processo.

\begin{figure}[t]
\center
\begin{tabular}{ll}
\textbf{Chamada de sistema}            & \textbf{Descrição}                                                                                                                           \\ \hline
int fork()                             & Cria um processo, retorna o PID do filho.                                                                                                    \\
int exit(int status)                   & Termina o processo atual; status retornado para wait(). Sem retorno.                                                                         \\
int wait(int *status)                  & \begin{tabular}[c]{@{}l@{}}Aguarda a finalização de um processo filho; status é colocado \\ em *status; retorna o PID do filho.\end{tabular} \\
int kill(int pid)                      & Encerra o processo PID. Retorna 0, ou -1 em caso de erro.                                                                                    \\
int getpid()                           & Retorna o PID do processo atual.                                                                                                             \\
int sleep(int n)                       & Pausa por n ciclos de clock.                                                                                                                 \\
int exec(char *file, char *argv{[}{]}) & \begin{tabular}[c]{@{}l@{}}Carrega um arquivo e executa com argumentos; só retorna em caso \\ de erro.\end{tabular}                          \\
char *sbrk(int n)                      & \begin{tabular}[c]{@{}l@{}}Expande a memória do processo em n bytes zero. Retorna o início da \\ nova memória.\end{tabular}                  \\
int open(char *file, int flags)        & \begin{tabular}[c]{@{}l@{}}Abre um arquivo; flags indicam leitura/escrita; retorna um fd \\ (descritor de arquivo).\end{tabular}             \\
int write(int fd, char *buf, int n)    & Escreve n bytes de buf no descritor fd; retorna n.                                                                                           \\
int read(int fd, char *buf, int n)     & Lê n bytes em buf; retorna quantidade lida, ou 0 se fim de arquivo.                                                                          \\
int close(int fd)                      & Libera o arquivo aberto fd.                                                                                                                  \\
int dup(int fd)                        & Retorna um novo descritor que referencia o mesmo arquivo que fd.                                                                             \\
int pipe(int p{[}{]})                  & Cria um pipe, insere descritores de leitura/escrita em p{[}0{]} e p{[}1{]}.                                                                  \\
int chdir(char *dir)                   & Altera o diretório atual.                                                                                                                    \\
int mkdir(char *dir)                   & Cria um novo diretório.                                                                                                                      \\
int mknod(char *file, int, int)        & Cria um arquivo de dispositivo.                                                                                                              \\
int fstat(int fd, struct stat *st)     & Insere informações de um arquivo aberto em *st.                                                                                              \\
int link(char *file1, char *file2)     & Cria outro nome (file2) para o arquivo file1.                                                                                                \\
int unlink(char *file)                 & Remove um arquivo.                                                                                                                          
\end{tabular}
\caption{Chamadas de sistema do xv6. Salvo indicação contrária, essas chamadas retornam 0 se não houver erro, e -1 em caso de erro.}
\label{fig:api}
\end{figure}

Um processo pode criar um novo processo usando a chamada de sistema
\indexcode{fork}.
\lstinline{fork}
cria no novo processo uma cópia exata da memória do processo que fez a chamada: ele copia as instruções, dados e
a pilha do processo original para o novo processo.
\lstinline{fork}
retorna em ambos os processos: o original e o novo.
No processo original, \lstinline{fork} retorna o PID do novo processo.
No processo novo, \lstinline{fork} retorna zero.
Os processos original e novo são frequentemente chamados de
\indextext{pai}
e
\indextext{filho}.

Por exemplo, considere o seguinte trecho de programa escrito na linguagem C~\cite{kernighan}:
\begin{lstlisting}[]
int pid = fork();
if(pid > 0){
  printf("pai: filho=%d\n", pid);
  pid = wait((int *) 0);
  printf("filho %d terminou\n", pid);
} else if(pid == 0){
  printf("filho: saindo\n");
  exit(0);
} else {
  printf("erro no fork\n");
}
\end{lstlisting}
A chamada de sistema
\indexcode{exit}
faz com que o processo que a chamou pare de executar e
libere recursos como memória e arquivos abertos.
Exit recebe um argumento inteiro de status,
convencionalmente 0 para indicar sucesso e 1 para indicar falha.
A chamada
\indexcode{wait}
retorna o PID de um filho finalizado (ou morto) do processo atual
e copia o status de saída do filho para o endereço passado para wait;
se nenhum dos filhos do chamador tiver finalizado,
\indexcode{wait}
aguarda que um finalize.
Se o processo pai não tiver filhos, \lstinline{wait} retorna imediatamente -1.
Se o pai não se importa com o status de saída do filho, pode
passar o endereço nulo (0) para
\lstinline{wait}.

No exemplo, as linhas de saída
\begin{lstlisting}[]
pai: filho=1234
filho: saindo
\end{lstlisting}
podem aparecer em qualquer ordem (ou até intercaladas), dependendo de quem
executar o \indexcode{printf} primeiro: o pai ou o filho.
Após o filho finalizar, o \indexcode{wait} do pai retorna, fazendo com que o pai imprima:
\begin{lstlisting}[]
pai: filho 1234 terminou
\end{lstlisting}
Embora o filho comece com o mesmo conteúdo de memória do pai,
pai e filho executam com memórias e registradores separados:
alterar uma variável em um não afeta o outro. Por exemplo, quando o valor de retorno de
\lstinline{wait}
é armazenado em
\lstinline{pid}
no processo pai,
isso não altera a variável
\lstinline{pid}
no filho. O valor de
\lstinline{pid}
no filho continuará sendo zero.

A chamada de sistema
\indexcode{exec}
substitui a memória do processo chamador por uma nova imagem
de memória carregada de um arquivo no sistema de arquivos.
O arquivo deve ter um formato específico, que indica quais partes contêm instruções,
dados, onde iniciar a execução, etc.
O xv6 usa o formato ELF, que será discutido com mais detalhes no Capítulo~\ref{CH:MEM}.
Normalmente, o arquivo é o resultado da compilação de um código-fonte.
Quando
\indexcode{exec}
é bem-sucedido, ele não retorna ao programa chamador;
em vez disso, as instruções carregadas começam a ser executadas a partir do ponto de entrada
indicado no cabeçalho ELF.
\lstinline{exec}
recebe dois argumentos: o nome do arquivo executável e um vetor de argumentos em formato string.
Por exemplo:
\begin{lstlisting}[]
char *argv[3];

argv[0] = "echo";
argv[1] = "hello";
argv[2] = 0;
exec("/bin/echo", argv);
printf("erro no exec\n");
\end{lstlisting}
Esse trecho substitui o programa chamador por uma instância
do programa 
\lstinline{/bin/echo}
executando com a lista de argumentos
\lstinline{echo hello}.
A maioria dos programas ignora o primeiro elemento do vetor de argumentos, que
convencionalmente é o nome do programa.

O shell do xv6 usa essas chamadas para executar programas em nome dos
usuários. A estrutura principal do shell é simples; veja
\lstinline{main}.
O loop principal lê uma linha de entrada do usuário com
\indexcode{getcmd}.
Depois, chama
\lstinline{fork},
que cria uma cópia do processo do shell.
O processo pai chama
\lstinline{wait},
enquanto o filho executa o comando. Por exemplo, se o usuário digitar
``\lstinline{echo hello}''
no shell,
\lstinline{runcmd}
será chamado com
``\lstinline{echo hello}''
como argumento.
\lstinline{runcmd}
executa o comando propriamente dito. Para
``\lstinline{echo hello}'',
ele chamaria
\lstinline{exec}.
Se
\lstinline{exec}
for bem-sucedido, o filho passará a executar as instruções de
\lstinline{echo}
em vez de
\lstinline{runcmd}.
Em algum momento,
\lstinline{echo}
chamará
\lstinline{exit},
fazendo com que o pai retorne de
\lstinline{wait}
em
\lstinline{main}.

Você pode se perguntar por que
\indexcode{fork}
e
\indexcode{exec}
não são combinadas em uma única chamada; veremos mais adiante que
o shell se aproveita dessa separação para implementar
redirecionamento de entrada/saída.
Para evitar o desperdício de criar um processo duplicado apenas para substituí-lo
imediatamente (com \lstinline{exec}),
os kernels modernos otimizam a implementação de
\lstinline{fork}
usando técnicas de memória virtual, como
copy-on-write (ver Seção~\ref{sec:pagefaults}).

O xv6 aloca a maior parte da memória do espaço do usuário
implicitamente:
\indexcode{fork}
aloca a memória necessária para a cópia da memória do pai no filho, e
\indexcode{exec}
aloca memória suficiente para carregar o executável.
Um processo que precisa de mais memória em tempo de execução (por exemplo, para
\indexcode{malloc})
pode chamar
\lstinline{sbrk(n)}
para expandir sua memória de dados em
\lstinline{n} bytes zero;
\indexcode{sbrk}
retorna o endereço da nova memória.


%% 
%% 	I/O and File descriptors
%% 
\section{E/S e Descritores de Arquivo}

Um 
\indextext{descritor de arquivo}
é um pequeno número inteiro que representa um objeto gerenciado pelo kernel
do qual um processo pode ler ou para o qual pode escrever.
Um processo pode obter um descritor de arquivo ao abrir um arquivo, diretório
ou dispositivo, ou ao criar um pipe, ou duplicando um descritor
existente.
Para simplificar, frequentemente nos referiremos ao objeto associado ao descritor de arquivo
simplesmente como um ``arquivo'';  
a interface de descritor de arquivo abstrai as diferenças entre
arquivos, pipes e dispositivos, fazendo com que todos pareçam fluxos de bytes.
Chamaremos leitura e escrita de \indextext{E/S} (entrada/saída).

Internamente, o kernel do xv6 usa o descritor de arquivo
como um índice em uma tabela por processo,
de modo que cada processo possui um espaço privado de descritores de arquivo,
começando pelo zero.
Por convenção, um processo lê do descritor 0 (entrada padrão),
escreve na saída padrão pelo descritor 1 e
escreve mensagens de erro no descritor 2 (erro padrão).
Como veremos, o shell aproveita essa convenção para implementar redirecionamento de E/S
e pipelines. O shell garante que sempre tenha três descritores de arquivo
abertos,
que, por padrão, estão conectados ao console.

As chamadas de sistema
\lstinline{read}
e
\lstinline{write}
leem bytes de e escrevem bytes em
arquivos abertos identificados por descritores de arquivo.
A chamada
\lstinline{read(fd, buf, n)}
lê no máximo
\lstinline{n}
bytes do descritor de arquivo
\lstinline{fd},
copia-os para
\lstinline{buf},
e retorna o número de bytes lidos.
Cada descritor de arquivo associado a um arquivo
possui um deslocamento associado.
\lstinline{read}
lê os dados a partir do deslocamento atual do arquivo e então avança
esse deslocamento pelo número de bytes lidos:
uma chamada subsequente a
\lstinline{read}
retornará os bytes seguintes aos da primeira chamada.
Quando não há mais bytes a serem lidos,
\lstinline{read}
retorna zero, indicando fim de arquivo.

A chamada
\lstinline{write(fd, buf, n)}
escreve
\lstinline{n}
bytes de
\lstinline{buf}
no descritor de arquivo
\lstinline{fd}
e retorna o número de bytes escritos.
Menos de
\lstinline{n}
bytes serão escritos apenas em caso de erro.
Assim como
\lstinline{read},
\lstinline{write}
escreve os dados no deslocamento atual e depois o avança
de acordo com o número de bytes escritos:
cada
\lstinline{write}
continua a partir de onde a anterior parou.

O trecho de programa a seguir (que representa a essência do programa
\lstinline{cat})
copia dados da entrada padrão
para a saída padrão. Se ocorrer algum erro, ele escreve uma mensagem
no erro padrão.
\begin{lstlisting}[]
char buf[512];
int n;

for(;;){
  n = read(0, buf, sizeof buf);
  if(n == 0)
    break;
  if(n < 0){
    fprintf(2, "erro de leitura\n");
    exit(1);
  }
  if(write(1, buf, n) != n){
    fprintf(2, "erro de escrita\n");
    exit(1);
  }
}
\end{lstlisting}
O importante a se notar nesse trecho é que o programa
\lstinline{cat}
não sabe se está lendo de um arquivo, do console ou de um pipe.
Da mesma forma,
\lstinline{cat}
não sabe se está escrevendo no console, em um arquivo, ou em outro lugar.
O uso de descritores de arquivo e a convenção de que o descritor 0
é a entrada e o descritor 1 é a saída permite uma implementação
simples do
\lstinline{cat}.

A chamada de sistema
\lstinline{close}
libera um descritor de arquivo, tornando-o disponível para reutilização por uma futura
chamada
\lstinline{open},
\lstinline{pipe}
ou
\lstinline{dup}
(ver abaixo).
Um novo descritor de arquivo alocado
é sempre o menor descritor não utilizado
pelo processo atual.

Descritores de arquivo e a chamada
\indexcode{fork}
interagem para tornar fácil a implementação de redirecionamento de E/S.
\lstinline{fork}
copia a tabela de descritores de arquivo do processo pai junto com sua memória,
de modo que o filho inicia com exatamente os mesmos arquivos abertos que o pai.
A chamada de sistema
\indexcode{exec}
substitui a memória do processo chamador, mas preserva sua tabela de arquivos.
Esse comportamento permite que o shell
implemente \indextext{redirecionamento de E/S} fazendo fork,
reabrindo os descritores de arquivo desejados no filho
e, então, chamando \lstinline{exec} para executar o novo programa.
Aqui está uma versão simplificada do código que um shell executa para o comando
\lstinline{cat < input.txt}:
\begin{lstlisting}[]
char *argv[2];

argv[0] = "cat";
argv[1] = 0;
if(fork() == 0) {
  close(0);
  open("input.txt", O_RDONLY);
  exec("cat", argv);
}
\end{lstlisting}
Após o processo filho fechar o descritor de arquivo 0,
\lstinline{open}
garantidamente usará esse descritor
para abrir
\lstinline{input.txt}:
0 será o menor descritor disponível.
O programa
\lstinline{cat}
executará com o descritor 0 (entrada padrão) referindo-se a
\lstinline{input.txt}.
Os descritores do processo pai não são afetados por essa
sequência, pois ela modifica apenas os descritores do filho.

O código de redirecionamento de E/S no shell do xv6 funciona exatamente dessa forma.
Lembre-se de que, nesse ponto do código, o shell já fez fork do
processo filho, e
\lstinline{runcmd}
chamará
\lstinline{exec}
para carregar o novo programa.

O segundo argumento de \lstinline{open} consiste em um conjunto de
flags, expressos como bits, que controlam o comportamento da chamada \lstinline{open}.
Os valores possíveis são definidos no cabeçalho de controle de arquivos (fcntl):
\lstinline{O_RDONLY},
\lstinline{O_WRONLY},
\lstinline{O_RDWR},
\lstinline{O_CREATE}, e
\lstinline{O_TRUNC},
que instruem \lstinline{open} a
abrir o arquivo para leitura,
ou para escrita,
ou para leitura e escrita,
a criar o arquivo se ele não existir,
e a truncar (zerar) o arquivo.

Agora fica claro por que é útil que
\lstinline{fork}
e
\lstinline{exec}
sejam chamadas separadas: entre elas, o shell tem a chance
de redirecionar a E/S do processo filho sem afetar a configuração de E/S do shell principal.
Pode-se imaginar uma hipotética chamada combinada
\lstinline{forkexec},
mas as opções para redirecionamento de E/S com essa chamada
seriam complicadas.
O shell teria que modificar sua própria E/S antes de chamar \lstinline{forkexec} (e então
desfazer as modificações); ou
\lstinline{forkexec} teria que aceitar instruções de redirecionamento como argumentos;
ou (menos desejável) cada programa como \lstinline{cat}
teria que ser modificado para fazer seu próprio redirecionamento.

Embora
\lstinline{fork}
copie a tabela de descritores de arquivos, o deslocamento interno de cada arquivo é compartilhado
entre pai e filho.
Considere este exemplo:
\begin{lstlisting}[]
if(fork() == 0) {
  write(1, "hello ", 6);
  exit(0);
} else {
  wait(0);
  write(1, "world\n", 6);
}
\end{lstlisting}
Ao final desse trecho, o arquivo associado ao descritor 1
conterá os dados
\lstinline{hello}
\lstinline{world}.
A chamada
\lstinline{write}
no processo pai
(que, graças ao
\lstinline{wait},
executa apenas após o filho terminar)
continua a partir do ponto onde a chamada
\lstinline{write}
do filho parou.
Esse comportamento ajuda a produzir saídas sequenciais de comandos do shell como:
\lstinline{(echo hello; echo world) >output.txt}.

A chamada
\lstinline{dup}
duplica um descritor de arquivo existente,
retornando um novo que se refere ao mesmo objeto de E/S.
Ambos os descritores compartilham o mesmo deslocamento, assim como os descritores
duplicados por
\lstinline{fork}.
Esta é outra maneira de escrever
\lstinline{hello world}
em um arquivo:
\begin{lstlisting}[]
fd = dup(1);
write(1, "hello ", 6);
write(fd, "world\n", 6);
\end{lstlisting}

Dois descritores compartilham o deslocamento se derivarem do
mesmo descritor original por meio de chamadas de
\lstinline{fork}
e
\lstinline{dup}.
Caso contrário, os descritores não compartilham deslocamento, mesmo que resultem de
chamadas
\lstinline{open}
para o mesmo arquivo.
\lstinline{dup}
permite que shells implementem comandos como:
\lstinline{ls existing-file non-existing-file >tmp1 2>&1}.
O trecho
\lstinline{2>&1}
informa ao shell que o descritor 2 deve ser uma duplicata do descritor 1.
Tanto o nome do arquivo existente quanto a mensagem de erro do arquivo
inexistente aparecerão no arquivo
\lstinline{tmp1}.
O shell do xv6 não suporta redirecionamento de E/S para o descritor de erro,
mas agora você sabe como implementá-lo.

Descritores de arquivos são uma abstração poderosa,
pois escondem os detalhes do que está conectado a eles:
um processo escrevendo no descritor 1 pode estar escrevendo em um
arquivo, em um dispositivo como o console ou em um pipe.

%% 
%% 	Pipes
%% 
\section{Pipes}

Um 
\indextext{pipe}
é um pequeno buffer no kernel exposto aos processos como um par de
descritores de arquivo: um para leitura e outro para escrita.
Escrever dados em uma extremidade do pipe
torna esses dados disponíveis para leitura na outra extremidade.
Pipes fornecem um meio de comunicação entre processos.

O código de exemplo a seguir executa o programa
\lstinline{wc}
com a entrada padrão conectada à extremidade de leitura de um pipe.
\begin{lstlisting}[]
int p[2];
char *argv[2];

argv[0] = "wc";
argv[1] = 0;

pipe(p);
if(fork() == 0) {
  close(0);
  dup(p[0]);
  close(p[0]);
  close(p[1]);
  exec("/bin/wc", argv);
} else {
  close(p[0]);
  write(p[1], "hello world\n", 12);
  close(p[1]);
}
\end{lstlisting}

O programa chama
\lstinline{pipe},
que cria um novo pipe e armazena os descritores de leitura e escrita
no array
\lstinline{p}.
Após o
\lstinline{fork},
tanto o processo pai quanto o processo filho possuem descritores de arquivo
referentes ao pipe.
O processo filho chama \lstinline{close} e \lstinline{dup} para fazer com que o descritor 0
passe a se referir à extremidade de leitura do pipe,
fecha os descritores do array
\lstinline{p},
e chama \lstinline{exec} para executar
\lstinline{wc}.
Quando
\lstinline{wc}
lê da entrada padrão, ele estará lendo do pipe.
O processo pai fecha o lado de leitura do pipe,
escreve no pipe,
e então fecha o lado de escrita.

Se não houver dados disponíveis, uma chamada
\lstinline{read}
em um pipe espera até que dados sejam escritos ou até que todos os
descritores referentes à extremidade de escrita sejam fechados;
neste último caso,
\lstinline{read}
retorna zero, assim como quando se chega ao final de um arquivo.
O fato de que
\lstinline{read}
bloqueia até que seja impossível novos dados chegarem
é uma das razões pelas quais é importante que o processo filho
feche a extremidade de escrita do pipe
antes de executar
\lstinline{wc}
no exemplo acima: se algum dos descritores do
\lstinline{wc}
ainda referisse a extremidade de escrita,
\lstinline{wc}
nunca veria o fim de arquivo.

O shell do xv6 implementa pipelines como
\lstinline{grep fork sh.c | wc -l}
de maneira semelhante ao código acima.
O processo filho cria um pipe para conectar a extremidade esquerda do pipeline
com a extremidade direita. Em seguida, ele chama
\lstinline{fork}
e
\lstinline{runcmd}
para a extremidade esquerda do pipeline,
e depois chama
\lstinline{fork}
e
\lstinline{runcmd}
para a extremidade direita, e espera que ambos terminem.
A extremidade direita do pipeline pode ser um comando que por sua vez inclua outro
pipe (por exemplo,
\lstinline{a | b | c}),
o que leva à criação de dois novos processos filhos (um para
\lstinline{b}
e outro para
\lstinline{c}).
Assim, o shell pode criar uma árvore de processos. As folhas dessa árvore são os comandos,
e os nós internos são processos que esperam até que seus filhos à esquerda e à direita terminem.

Pipes podem parecer não ser mais poderosos do que arquivos temporários:
o pipeline
\begin{lstlisting}[]
echo hello world | wc
\end{lstlisting}
poderia ser implementado sem pipes como:
\begin{lstlisting}[]
echo hello world >/tmp/xyz; wc </tmp/xyz
\end{lstlisting}
Mas os pipes possuem ao menos três vantagens sobre arquivos temporários
nesta situação.

Primeiro, pipes se limpam automaticamente;
com redirecionamento de arquivos, o shell precisaria
se preocupar em remover
\lstinline{/tmp/xyz}
ao final da execução.

Segundo, pipes podem transmitir fluxos de dados arbitrariamente longos,
enquanto o redirecionamento para arquivos exige espaço livre em disco
suficiente para armazenar todos os dados.

Terceiro, pipes permitem a execução paralela das etapas do pipeline,
enquanto a abordagem com arquivos exige que o primeiro programa finalize
antes do segundo começar.

% Quarto, se você estiver implementando comunicação entre processos,
% pipes com leituras e escritas bloqueantes são mais eficientes
% do que arquivos com semântica não bloqueante.

%% 
%% 	File system
%% 
\section{Sistema de arquivos}

O sistema de arquivos do xv6 fornece arquivos de dados,
que contêm arrays de bytes não interpretados,
e diretórios, que
contêm referências nomeadas para arquivos de dados e outros diretórios.
Os diretórios formam uma árvore, começando
em um diretório especial chamado
\indextext{raiz} (root).
Um
\indextext{caminho} (path)
como
\lstinline{/a/b/c}
refere-se ao arquivo ou diretório chamado
\lstinline{c}
dentro do diretório chamado
\lstinline{b},
dentro do diretório chamado
\lstinline{a},
no diretório raiz
\lstinline{/}.
Caminhos que não começam com
\lstinline{/}
são avaliados em relação ao
\indextext{diretório atual}
do processo chamador,
que pode ser alterado com a chamada de sistema
\lstinline{chdir}.
Ambos os trechos de código a seguir abrem o mesmo arquivo
(assumindo que todos os diretórios existam):

\begin{lstlisting}[]
chdir("/a");
chdir("b");
open("c", O_RDONLY);

open("/a/b/c", O_RDONLY);
\end{lstlisting}

O primeiro trecho altera o diretório atual do processo para
\lstinline{/a/b};
o segundo não altera nem depende do diretório atual do processo.

Existem chamadas de sistema para criar novos arquivos e diretórios:
\lstinline{mkdir}
cria um novo diretório,
\lstinline{open}
com a flag
\lstinline{O_CREATE}
cria um novo arquivo de dados,
e
\lstinline{mknod}
cria um novo arquivo de dispositivo.
Este exemplo ilustra os três casos:
\begin{lstlisting}[]
mkdir("/dir");
fd = open("/dir/file", O_CREATE|O_WRONLY);
close(fd);
mknod("/console", 1, 1);
\end{lstlisting}

\lstinline{mknod}
cria um arquivo especial que se refere a um dispositivo.
Associados a um arquivo de dispositivo estão
o número major e o número minor do dispositivo
(os dois argumentos de
\lstinline{mknod}),
que identificam de forma única um dispositivo do kernel.
Quando um processo posteriormente abre um arquivo de dispositivo, o kernel
desvia as chamadas de sistema
\lstinline{read}
e
\lstinline{write}
para a implementação do dispositivo no kernel,
em vez de passá-las ao sistema de arquivos.

O nome de um arquivo é distinto do próprio arquivo;
o mesmo arquivo subjacente, chamado
\indextext{inode},
pode ter vários nomes,
chamados
\indextext{links}.
Cada link consiste em uma entrada de diretório;
essa entrada contém um nome de arquivo e uma referência
a um inode.
Um inode armazena
\indextext{metadados}
sobre um arquivo, incluindo
seu tipo (arquivo, diretório ou dispositivo),
seu tamanho,
a localização do conteúdo do arquivo no disco
e o número de links associados ao arquivo.

A chamada de sistema
\lstinline{fstat}
recupera informações do inode ao qual um
descritor de arquivo se refere.
Ela preenche uma
\lstinline{struct stat},
definida em
\lstinline{stat.h} \fileref{kernel/stat.h}
como:
\begin{lstlisting}[]
#define T_DIR     1   // Diretório
#define T_FILE    2   // Arquivo
#define T_DEVICE  3   // Dispositivo

struct stat {
  int dev;     // Dispositivo de disco do sistema de arquivos
  uint ino;    // Número do inode
  short type;  // Tipo de arquivo
  short nlink; // Número de links para o arquivo
  uint64 size; // Tamanho do arquivo em bytes
};
\end{lstlisting}

A chamada de sistema
\lstinline{link}
cria outro nome no sistema de arquivos
referente ao mesmo inode de um arquivo existente.
Este trecho cria um novo arquivo com os nomes
\lstinline{a}
e
\lstinline{b}:
\begin{lstlisting}[]
open("a", O_CREATE|O_WRONLY);
link("a", "b");
\end{lstlisting}
Ler ou escrever em
\lstinline{a}
é o mesmo que ler ou escrever em
\lstinline{b}.
Cada inode é identificado por um número único de
\textit{inode}.
Após a execução do código acima, é possível
verificar que
\lstinline{a}
e
\lstinline{b}
referem-se ao mesmo conteúdo subjacente, inspecionando o resultado de
\lstinline{fstat}:
ambos retornarão o mesmo número de inode
(\lstinline{ino}),
e a contagem de
\lstinline{nlink}
será 2.

A chamada de sistema
\lstinline{unlink}
remove um nome do sistema de arquivos.
O inode e o espaço de disco associado ao conteúdo do arquivo
só são liberados quando a contagem de links for zero e
nenhum descritor de arquivo estiver referenciando o arquivo.
Assim, ao adicionar:
\begin{lstlisting}[]
unlink("a");
\end{lstlisting}
ao código anterior, o inode
e o conteúdo do arquivo continuam acessíveis por meio de
\lstinline{b}.
Além disso, o trecho:
\begin{lstlisting}[]
fd = open("/tmp/xyz", O_CREATE|O_RDWR);
unlink("/tmp/xyz");
\end{lstlisting}
é uma forma idiomática de criar um inode temporário
sem nome,
que será automaticamente removido quando o processo fechar
\lstinline{fd}
ou terminar sua execução.

O Unix fornece
utilitários de arquivos que podem ser chamados a partir do shell
como programas de usuário, por exemplo
\lstinline{mkdir},
\lstinline{ln}
e
\lstinline{rm}.
Esse design permite que qualquer pessoa estenda a interface de linha de comando
adicionando novos programas de usuário.
Em retrospecto, esse plano parece óbvio,
mas outros sistemas projetados na mesma época que o Unix muitas vezes embutiam tais comandos diretamente
no shell (e até mesmo o shell no kernel).

Uma exceção é o comando
\lstinline{cd},
que é incorporado ao próprio shell.
\lstinline{cd}
precisa alterar o diretório de trabalho atual do
próprio shell.
Se
\lstinline{cd}
fosse executado como um comando comum, o shell faria um
\textit{fork} de um processo filho, e o processo filho executaria
\lstinline{cd},
mudando o diretório de trabalho do
\textit{filho},
mas o diretório de trabalho do processo pai (ou seja, o shell)
permaneceria inalterado.

%% 
%% 	Real world
%% 
\section{Mundo real}

A combinação do Unix de descritores de arquivos “padrão”,
pipes e uma sintaxe de shell conveniente para operações sobre eles
foi um grande avanço na escrita de programas reutilizáveis de propósito geral.
Essa ideia deu origem a uma cultura de “ferramentas de software” que foi
responsável por grande parte do poder e da popularidade do Unix,
e o shell foi a primeira chamada “linguagem de script”.
A interface de chamadas de sistema do Unix persiste até hoje em sistemas como
BSD, Linux e macOS.

A interface de chamadas de sistema do Unix foi padronizada através do padrão
Portable Operating System Interface (POSIX).
O xv6
\textit{não}
é compatível com POSIX: ele carece de muitas chamadas de sistema (incluindo algumas básicas como
\lstinline{lseek}),
e muitas das chamadas que ele fornece diferem do padrão.
Nossos principais objetivos com o xv6 são
simplicidade e clareza, ao mesmo tempo em que fornece uma interface de chamadas de sistema simples e semelhante ao UNIX.
Várias pessoas estenderam o xv6 com algumas chamadas de sistema adicionais e uma biblioteca C simples para executar programas Unix básicos.
No entanto, kernels modernos fornecem muitas mais chamadas de sistema e tipos de serviços do kernel do que o xv6.
Por exemplo, eles suportam redes, sistemas de janelas, threads em nível de usuário,
drivers para muitos dispositivos e assim por diante.
Kernels modernos evoluem continuamente e rapidamente,
e oferecem muitos recursos além do POSIX.

O Unix unificou o acesso a múltiplos tipos de recursos (arquivos,
diretórios e dispositivos) com um único conjunto de interfaces
baseadas em nomes de arquivos e descritores de arquivos.
Essa ideia pode ser estendida a mais tipos de recursos;
um bom exemplo é o sistema Plan 9~\cite{Presotto91plan9},
que aplicou o conceito de “recursos são arquivos”
a redes, gráficos e mais.
No entanto, a maioria dos sistemas operacionais derivados do Unix
não seguiu esse caminho.

O sistema de arquivos e os descritores de arquivos têm sido abstrações poderosas.
Ainda assim, existem outros modelos para interfaces de sistemas operacionais.
O Multics, predecessor do Unix,
abstraía o armazenamento de arquivos de forma que o tornava semelhante à memória,
resultando em uma interface com uma abordagem bastante diferente.
A complexidade do projeto do Multics teve influência direta
nos projetistas do Unix, que buscaram criar algo mais simples.

O xv6 não fornece a noção de usuários ou proteção
entre usuários; em termos de Unix, todos os processos no xv6
são executados como root.

Este livro examina como o xv6 implementa sua interface semelhante ao Unix,
mas as ideias e conceitos se aplicam a mais do que apenas ao Unix.
Qualquer sistema operacional deve multiplexar processos sobre o
hardware subjacente, isolar processos entre si
e fornecer mecanismos para comunicação interprocessos controlada.
Após estudar o xv6, você deverá ser capaz de
analisar outros sistemas operacionais mais complexos
e reconhecer neles os conceitos fundamentais do xv6.

%%
\section{Exercícios}
%%

\begin{enumerate}

\item Escreva um programa que utilize chamadas de sistema UNIX para
fazer um “ping-pong” de um byte entre dois processos usando um par
de pipes, um para cada direção. Meça o
desempenho do programa, em trocas por segundo.

\end{enumerate}

